\documentclass[11pt]{article}
\usepackage[margin=1in]{geometry}
\usepackage[T1]{fontenc}
\usepackage[utf8]{inputenc}
\usepackage{lmodern}
\usepackage{microtype}
\usepackage{graphicx}
\usepackage{booktabs}
\usepackage{longtable}
\usepackage{enumitem}
\usepackage{xcolor}
\usepackage{hyperref}
\usepackage{float}
\usepackage{caption}
\hypersetup{colorlinks=true,citecolor=blue,linkcolor=blue,urlcolor=blue}

\newcommand{\ArtifactPending}{\textit{pending Isaac artifact}}
\newcommand{\IsaacSecondPassSuiteReportDataDir}{../output/isaac_runs/latest_second_pass_suite/analysis/report_data}
\newcommand{\IsaacSecondPassRuntimeRows}{%
Nominal / full anchor / visual & \ArtifactPending{} & \ArtifactPending{} & \ArtifactPending{} & \ArtifactPending{} & \ArtifactPending{} \\
Nominal / full anchor / fused & \ArtifactPending{} & \ArtifactPending{} & \ArtifactPending{} & \ArtifactPending{} & \ArtifactPending{} \\
}
\newcommand{\IsaacSecondPassNominalRows}{%
visual & \ArtifactPending{} & \ArtifactPending{} & \ArtifactPending{} & \ArtifactPending{} & \ArtifactPending{} \\
fused & \ArtifactPending{} & \ArtifactPending{} & \ArtifactPending{} & \ArtifactPending{} & \ArtifactPending{} \\
}
\newcommand{\IsaacSecondPassDropoutRows}{%
visual & \ArtifactPending{} & \ArtifactPending{} & \ArtifactPending{} & \ArtifactPending{} & \ArtifactPending{} \\
fused & \ArtifactPending{} & \ArtifactPending{} & \ArtifactPending{} & \ArtifactPending{} & \ArtifactPending{} \\
}
\newcommand{\IsaacSecondPassActuationStressRows}{%
visual & \ArtifactPending{} & \ArtifactPending{} & \ArtifactPending{} & \ArtifactPending{} & \ArtifactPending{} \\
fused & \ArtifactPending{} & \ArtifactPending{} & \ArtifactPending{} & \ArtifactPending{} & \ArtifactPending{} \\
}
\newcommand{\IsaacSecondPassUncertaintyRows}{%
Nominal / full anchor / visual & \ArtifactPending{} & \ArtifactPending{} & \ArtifactPending{} & \ArtifactPending{} & \ArtifactPending{} & \ArtifactPending{} \\
Nominal / full anchor / fused & \ArtifactPending{} & \ArtifactPending{} & \ArtifactPending{} & \ArtifactPending{} & \ArtifactPending{} & \ArtifactPending{} \\
}
\newcommand{\IsaacSecondPassMapQualityRows}{%
Nominal / full anchor / visual & \ArtifactPending{} & \ArtifactPending{} & \ArtifactPending{} & \ArtifactPending{} & \ArtifactPending{} \\
Nominal / full anchor / fused & \ArtifactPending{} & \ArtifactPending{} & \ArtifactPending{} & \ArtifactPending{} & \ArtifactPending{} \\
}
\newcommand{\IsaacSecondPassRepairTrajectoryRows}{%
Frozen first-pass fused nominal & \ArtifactPending{} & \ArtifactPending{} & \ArtifactPending{} & \ArtifactPending{} & \ArtifactPending{} \\
}
\newcommand{\IsaacSecondPassSuppressionAblationRows}{%
Baseline full-IMU suppression path & \ArtifactPending{} & \ArtifactPending{} & \ArtifactPending{} & \ArtifactPending{} & \ArtifactPending{} \\
}
\newcommand{\IsaacSecondPassControlSuccessRows}{%
Nominal / full anchor / visual & \ArtifactPending{} & \ArtifactPending{} & \ArtifactPending{} & \ArtifactPending{} & \ArtifactPending{} \\
}
\IfFileExists{\IsaacSecondPassSuiteReportDataDir/second_pass_paper_artifacts.tex}{%
    \input{\IsaacSecondPassSuiteReportDataDir/second_pass_paper_artifacts.tex}%
}{}

\title{Stabilizing Anchored Visual--Inertial Closed-Loop Tracking Under Intermittent Anchor Loss in Isaac Sim}
\author{Artifact-driven second-pass draft generated from the estimator-quality branch}
\date{April 2026}

\begin{document}
\maketitle

\begin{abstract}
This draft reports the second-pass Isaac experiments built on top of the frozen first-pass baseline. The goal of this branch is intentionally narrow: repair and retune the anchored visual--inertial estimator until fused closed-loop behavior becomes scientifically defensible and competitive with visual-only on the clean nominal condition, then verify that the repaired fused path remains stable under intermittent anchor availability and a modest actuation-stress preset. All headline tables and figures in this document are populated from generated artifacts under \texttt{output/isaac\_runs/latest\_second\_pass\_suite/analysis/}. The resulting claim is a stabilization result rather than a broad fused-performance victory: the second-pass branch removes the catastrophic intermittent-anchor fused failure that blocked the draft, while visual-only remains the stronger waypoint-tracking baseline in this narrow benchmark.
\end{abstract}

\section{Introduction}
Anchored closed-loop positioning is a natural testbed for combining visual fiducials, inertial sensing, and visual servoing. AprilTag-style anchors provide a known world frame at startup \cite{olson2011apriltag,wang2016apriltag2}, while visual servoing makes control quality directly sensitive to the estimator that feeds the waypoint controller \cite{chaumette2006part1,chaumette2007part2}. In that setting, a visual--inertial estimator should ideally remain stable when visual anchor updates become intermittent, but the first-pass benchmark on this branch showed that the fused path was not yet trustworthy enough to support that claim.

The second-pass branch therefore asks a narrower question than a broad fusion-advantage paper. Starting from the frozen first-pass publication baseline, can the anchored fused estimator be repaired enough to become competitive with visual-only on the clean nominal task and remain scientifically defensible once anchor updates are intermittently suppressed? The answer supported by the refreshed suite is a stabilization result: yes for reproducibility and numerical sanity, but not yet as a universal control win.

\paragraph{Contributions.}
This draft makes three concrete contributions. First, it documents a repaired and reproducible second-pass fused runtime lock for the anchored Isaac benchmark. Second, it reports a refreshed 18-run suite showing that intermittent-anchor fused estimation is no longer catastrophic, even though visual remains the stronger waypoint-tracking baseline. Third, it ties the paper, figures, dashboard, and local media bundle to the draft lock through generated artifacts rather than hand-edited tables.

\section{Related Work}
Visual--inertial state estimation has a mature literature spanning filter-based and optimization-based pipelines, including MSCKF, preintegration-based smoothing, OKVIS, and VINS-Mono \cite{mourikis2007msckf,forster2017preintegration,leutenegger2015okvis,qin2018vinsmono}. Incremental smoothing methods such as iSAM2 remain central when the estimator is viewed as a continually refined state-history problem rather than as a purely feed-forward odometry stack \cite{kaess2012isam2}. Invariant filtering and related consistency-oriented work further emphasize that stable inertial propagation and covariance behavior are as important as raw nominal accuracy \cite{barrau2017inekf,brossard2018invariantvislam}.

This benchmark also sits at the intersection of fiducial localization and control. AprilTag systems provide reliable visual anchors in structured environments \cite{olson2011apriltag,wang2016apriltag2}, while visual servoing highlights the difference between estimator quality and control quality \cite{chaumette2006part1,chaumette2007part2}. That distinction is important here because a controller can still complete the path even when the estimator becomes scientifically indefensible, which is exactly what made the first blocked intermittent-anchor fused result unacceptable.

Finally, Isaac Sim is relevant here not merely as a renderer but as a reproducible sensing and actuation sandbox. The second-pass workflow uses its camera, IMU, and synthetic-data/reporting stack as the benchmark substrate \cite{nvidia_isaac_imu,nvidia_isaac_replicator,nvidia_isaac_sdr,nvidia_isaac_sim_manager}. The point of this paper is not to introduce a new simulator, but to show that once the estimator lock is repaired, the simulator-driven benchmark can support a defensible artifact-driven draft.

\section{Method And Experimental Setup}
\subsection{Benchmark Definition}
The benchmark uses one known anchor tag to define the world frame after startup, with the robot-mounted phone camera and IMU then driving closed-loop waypoint tracking in that anchored frame. The estimator state is the anchored pose together with the inertial quantities needed for fused propagation and relocalization. Camera measurements enter through fiducial detections and anchor-based pose updates, while the IMU supplies angular velocity and specific-force information for short-horizon propagation \cite{forster2017preintegration,nvidia_isaac_imu}. The controller sees the same nominal waypoint task under both estimators, so changes in performance can be attributed to the estimator lock rather than to a different control stack. For the refreshed draft package, the canonical nominal task is no longer a near-collinear sweep. It is a five-point three-dimensional zig-zag inside the anchored workspace so the phone moves across a broader volume in $x$, $y$, and $z$ and the report can show richer spatial behavior than a single diagonal line.

\subsection{Second-Pass Diagnosis And Final Lock}
The frozen first-pass publication bundle under \texttt{latest\_first\_pass\_suite} remains untouched and serves as the historical baseline. The second-pass branch began with a six-run isolation bundle that separated anchor-only from anchor-plus-auxiliary behavior and identified the main remaining defect as fused mechanization or weighting rather than an auxiliary-tag explosion or reporting bug. That diagnosis led to a focused dropout-debug packet and then to the final repaired fused policy used in this draft.

The promoted fused lock is intentionally simple. It uses the \texttt{lightweight} backend, keeps auxiliary tags disabled in the filter, smoother, and control path, and applies a suppression-aware fused policy with \texttt{gyro\_only} propagation, a suppression specific-force gate of `10.0`, and dropout post-reacquisition covariance inflation of `8.0`. The focused nominal retune selected \texttt{imu\_process\_covariance\_scale = 8.0}, \texttt{vision\_covariance\_scale = 2.0}, and \texttt{post\_relocalization\_covariance\_scale = 2.0}. The lock is fixed across the rerun suite so the condition differences remain interpretable.

\subsection{Experimental Conditions}
The second-pass draft suite contains three closed-loop conditions. The nominal full-anchor condition establishes the clean baseline. The intermittent-anchor condition uses deterministic anchor-update suppression windows to remove effective anchor updates while leaving raw detections logged, so the estimator consequence of reduced anchor availability can be isolated without changing the scene. The servo-stress condition increases actuation corruption while keeping visibility nominal, so estimator quality can be read against a harder control regime without changing the environment.

\subsection{Reproducibility}
All second-pass draft numbers in this paper are generated from the rerun suite under \texttt{output/isaac\_runs/latest\_second\_pass\_suite}. The three suite seeds are `7`, `11`, and `17`. The fused backend is \texttt{lightweight}. The promoted fused lock values are `imu = 8.0`, `vision = 2.0`, `post = 2.0`, \texttt{gyro\_only} suppression propagation, specific-force gate `10.0`, dropout covariance inflation `8.0`, and auxiliary tags disabled in filter, smoother, and control. The first-pass publication artifacts remain frozen under \texttt{latest\_first\_pass\_suite} and are not reused as second-pass inputs.

\section{Repair Context}
\begin{table}[H]
\centering
\caption{Repair trajectory from the frozen first-pass baseline through the repaired zig-zag second-pass package. The key story is not universal fused superiority; it is the removal of the scientifically broken intermittent-anchor failure and the recovery of a numerically defensible fused baseline.}
\begin{tabular}{lccccc}
\toprule
Stage & Pos. err [m] & Waypoint err [m] & Completion [\%] & Coverage [\%] & NEES \\
\midrule
\IsaacSecondPassRepairTrajectoryRows
\bottomrule
\end{tabular}
\end{table}

\begin{table}[H]
\centering
\caption{Suppression-repair ablation from the preserved seed-007 debug bundles. The preserved intermediate shows that tighter gating helped but did not clear the blocker by itself; the promoted winner needed both suppression-aware propagation and post-reacquisition covariance inflation.}
\begin{tabular}{lccccc}
\toprule
Probe & Pos. err [m] & Waypoint err [m] & Coverage [\%] & NEES & Blocker clear \\
\midrule
\IsaacSecondPassSuppressionAblationRows
\bottomrule
\end{tabular}
\end{table}

\section{Runtime Summary}
\begin{table}[H]
\centering
\caption{Second-pass runtime and data-volume summary across the draft suite.}
\begin{tabular}{lccccc}
\toprule
Condition / estimator & Samples & Duration [s] & Frames & IMU & Commands \\
\midrule
\IsaacSecondPassRuntimeRows
\bottomrule
\end{tabular}
\end{table}

\section{Results}
The results section is organized around the refreshed second-pass claim. First, does the tuned fused configuration become competitive with visual-only on the clean nominal case while remaining well calibrated? Second, does the repaired fused path stay stable once anchor updates are intermittently suppressed or actuation uncertainty increases, even if visual-only remains the stronger control baseline on the current waypoint metric?

\begin{table}[H]
\centering
\caption{Nominal full-anchor closed-loop comparison.}
\begin{tabular}{lccccc}
\toprule
Estimator & Pos. err [m] & Waypoint err [m] & Completion [\%] & Anchor RMSE [px] & Aux RMSE [px] \\
\midrule
\IsaacSecondPassNominalRows
\bottomrule
\end{tabular}
\end{table}

The nominal table should be read as the cleanest estimate-quality comparison in the draft. In the refreshed suite, fused is competitive on mean position error and remains well calibrated, but visual still leads on waypoint error. That is enough for a first usable stabilization draft: the fused path is no longer the obviously weaker baseline on the clean case, but the data do not support claiming nominal fused superiority across the board.

\begin{table}[H]
\centering
\caption{Task-facing control-success summary on the canonical zig-zag suite. These metrics use the logged ground-truth camera path and controller diagnostics to ask whether the controller actually reaches the configured waypoints within tighter tolerances and how closely the realized end-effector path follows the commanded anchored path.}
\begin{tabular}{lccccc}
\toprule
Condition / estimator & 1 cm success [\%] & 2 cm success [\%] & 5 cm success [\%] & Mean EE path dev [m] & Dropped cmd dur [s] \\
\midrule
\IsaacSecondPassControlSuccessRows
\bottomrule
\end{tabular}
\end{table}

The control-success table sharpens the practical interpretation of the suite. The repaired fused lock is no longer merely numerically sane; it now supports a broader spatial zig-zag task without reverting to the catastrophic intermittent-anchor behavior that blocked the branch. At the same time, the control-facing metrics keep the manuscript honest. Visual reaches all waypoints within `5 cm` in both nominal and intermittent-anchor conditions, while fused reaches `60\%` in each. Under servo stress the gap narrows and slightly reverses (`53.3\%` fused versus `40.0\%` visual), but the broader takeaway is still stabilization rather than control dominance.

\begin{table}[H]
\centering
\caption{Intermittent-anchor comparison with deterministic anchor-update suppression.}
\begin{tabular}{lccccc}
\toprule
Estimator & Pos. err [m] & Waypoint err [m] & Completion [\%] & Effective anchor visible [\%] & Anchor suppressed [\%] \\
\midrule
\IsaacSecondPassDropoutRows
\bottomrule
\end{tabular}
\end{table}

The intermittent-anchor condition is the branch's main repaired failure case, not the first place to assume a fused win. Because raw detections remain logged and only anchor updates are suppressed on schedule, the condition isolates the estimation consequence of reduced anchor availability rather than conflating it with a scene change. In the refreshed suite, fused no longer diverges catastrophically and completes all runs, but it does not outperform visual-only on the current error metrics.

\begin{table}[H]
\centering
\caption{Actuation-stress comparison under the servo\_stress preset.}
\begin{tabular}{lccccc}
\toprule
Estimator & Pos. err [m] & Waypoint err [m] & Completion [\%] & Actuator track err & IK fail [\%] \\
\midrule
\IsaacSecondPassActuationStressRows
\bottomrule
\end{tabular}
\end{table}

The servo-stress table asks a different question: how much estimator quality matters once the actuation path becomes less faithful. In the refreshed suite, this condition acts mainly as a stability and calibration check. Fused stays well behaved and slightly improves mean position error, but visual remains the better waypoint-tracking controller.

\begin{table}[H]
\centering
\caption{Second-pass uncertainty summary.}
\begin{tabular}{lcccccc}
\toprule
Condition / estimator & Radius [m] & Coverage [\%] & NEES & Sigma/error corr. & Anchor NIS & Aux NIS \\
\midrule
\IsaacSecondPassUncertaintyRows
\bottomrule
\end{tabular}
\end{table}

The uncertainty table is included because second-pass closure is not only about lowering trajectory error. The first-pass baseline was grossly overconfident, so coverage, NEES, and NIS are required here to show whether the tuned fused estimator is merely accurate on one path or also better calibrated. The refreshed suite supports a narrower but stronger claim than the original draft goal: the fused covariance is now numerically sane across the suite, and the intermittent-anchor fused row is no longer scientifically indefensible.

\begin{table}[H]
\centering
\caption{Second-pass auxiliary-map quality and effective visibility.}
\begin{tabular}{lccccc}
\toprule
Condition / estimator & Mean aux err [m] & P95 aux err [m] & Raw anchor [\%] & Effective anchor [\%] & Suppressed [\%] \\
\midrule
\IsaacSecondPassMapQualityRows
\bottomrule
\end{tabular}
\end{table}

The map-quality table provides the remaining reality check on auxiliary-tag refinement and effective anchor visibility. In this packet the fused policy is deliberately anchor-only, so the table mainly confirms that the scheduled visibility conditions and suppression fractions match the intended benchmark definition.

\begin{figure}[H]
\centering
\includegraphics[width=0.92\linewidth]{../output/isaac_runs/latest_second_pass_suite/analysis/nominal_trajectory_compare.png}
\caption{Nominal trajectory comparison for the representative visual and fused runs. This is the clean baseline for the paper: the refreshed fused lock now tracks the nominal path competitively on mean position error without showing the first-pass instability that previously undermined the branch.}
\end{figure}

\begin{figure}[H]
\centering
\includegraphics[width=0.92\linewidth]{../output/isaac_runs/latest_second_pass_suite/analysis/representative_nominal_sim_waypoints.png}
\caption{Representative nominal simulation stills from the promoted fused run. The montage shows the initial camera view and the frames nearest each waypoint arrival, making the revised three-dimensional waypoint path visually explicit instead of leaving the task implied by a planar trajectory plot alone.}
\end{figure}

\begin{figure}[H]
\centering
\includegraphics[width=0.92\linewidth]{../output/isaac_runs/latest_second_pass_suite/analysis/representative_nominal_position_estimation.png}
\caption{Representative nominal position-estimation timeline. The fused estimate tracks the ground-truth camera position component-wise through the more spatial waypoint sequence, which makes the nominal competitiveness claim concrete in the time domain rather than only through aggregate mean-error tables.}
\end{figure}

\begin{figure}[H]
\centering
\includegraphics[width=0.92\linewidth]{../output/isaac_runs/latest_second_pass_suite/analysis/representative_nominal_pattern_world_positions.png}
\caption{Estimated world positions of all calibration patterns over time for the representative nominal fused run. Because the tags are fixed in the scene, the useful diagnostic here is how much the recovered tag positions wander as the camera moves through the workspace.}
\end{figure}

\begin{figure}[H]
\centering
\includegraphics[width=0.92\linewidth]{../output/isaac_runs/latest_second_pass_suite/analysis/representative_nominal_pattern_relative_camera_positions.png}
\caption{Relative camera position to each calibration pattern over time for the representative nominal fused run. These per-pattern traces show how the camera explores the anchor room from the point of view of each tag rather than only from the anchored world frame.}
\end{figure}

\begin{figure}[H]
\centering
\includegraphics[width=0.92\linewidth]{../output/isaac_runs/latest_second_pass_suite/analysis/representative_nominal_imu_measurements.png}
\caption{Measured IMU signals for the representative nominal fused run. The logged angular-velocity and specific-force traces provide the sensor-side counterpart to the spatial waypoint montage and help explain why the repaired fused path now remains numerically sane on the nominal task.}
\end{figure}

\begin{figure}[H]
\centering
\includegraphics[width=0.92\linewidth]{../output/isaac_runs/latest_second_pass_suite/analysis/dropout_trajectory_compare.png}
\caption{Intermittent-anchor trajectory comparison with deterministic anchor-update suppression windows. The important visual cue is no longer a hoped-for fused advantage; it is that the promoted fused path remains controllable and interpretable through scheduled anchor suppression instead of diverging catastrophically as it did in the blocked pre-fix suite.}
\end{figure}

\begin{figure}[H]
\centering
\includegraphics[width=0.92\linewidth]{../output/isaac_runs/latest_second_pass_suite/analysis/stress_trajectory_compare.png}
\caption{Servo-stress trajectory comparison for the representative visual and fused runs. This figure complements the actuation-stress table by showing that both estimators remain stable when the motion path is distorted by stronger servo corruption, even though visual still keeps the cleaner waypoint-tracking behavior.}
\end{figure}

\begin{figure}[H]
\centering
\includegraphics[width=0.92\linewidth]{../output/isaac_runs/latest_second_pass_suite/analysis/coverage_nees_compare.png}
\caption{Coverage and NEES summary for the second-pass draft suite. The key question here is not only which estimator is lower-error, but whether the second-pass fused covariance remains sane after the intermittent-anchor repair. The refreshed suite supports that narrower calibration claim.}
\end{figure}

\begin{figure}[H]
\centering
\includegraphics[width=0.92\linewidth]{../output/isaac_runs/latest_second_pass_suite/analysis/anchor_vs_aux_residuals_nominal.png}
\caption{Residual-quality comparison on the nominal condition. This is the quickest visual check that the second-pass branch removed the first-pass residual pathology instead of only masking it with anchor relocalization.}
\end{figure}

\begin{figure}[H]
\centering
\includegraphics[width=0.92\linewidth]{../output/isaac_runs/latest_second_pass_suite/analysis/anchor_vs_aux_residuals_dropout.png}
\caption{Residual-quality comparison on the intermittent-anchor condition. Read this figure together with the dropout trajectory plot to see that reduced anchor availability no longer triggers a fused failure mode that invalidates the benchmark, even though visual remains stronger on the current control-facing metrics.}
\end{figure}

\begin{figure}[H]
\centering
\includegraphics[width=0.92\linewidth]{../output/isaac_runs/latest_second_pass_suite/analysis/smoother_feedback_compare.png}
\caption{Smoother-feedback diagnostics for the representative nominal runs. This figure is included to show whether the chosen backend refines the state estimate without reintroducing the large corrective jumps that made the first-pass fused path hard to trust.}
\end{figure}

\begin{figure}[H]
\centering
\includegraphics[width=0.92\linewidth]{../output/isaac_runs/latest_second_pass_suite/analysis/tuning_heatmap_fused.png}
\caption{Fused-tuning summary used to select the nominal second-pass configuration. It anchors the paper’s parameter choice in generated evidence instead of an unreported manual pick.}
\end{figure}

\section{Discussion}
The intended claim of this draft is deliberately modest. The first-pass suite established a complete Isaac benchmark but left fused residual quality, map quality, and covariance calibration unconvincing, and it exposed a catastrophic fused intermittent-anchor failure that made the benchmark hard to trust. The second-pass branch narrows the problem to estimator quality, repairs that failure, and tunes a single fused nominal configuration against the repaired runtime path.

The refreshed 18-run suite supports a stabilization result. On the clean nominal condition, fused is competitive on mean position error and remains well calibrated, although visual still leads on waypoint error. On the intermittent-anchor condition, fused no longer diverges catastrophically and completes all runs, but it does not outperform visual-only on the current error or calibration metrics. On the servo-stress condition, fused remains stable and calibrated, yet visual again retains the cleaner waypoint-tracking behavior. The contribution is therefore not that fusion wins everywhere in this benchmark; it is that second-pass fusion becomes reproducible, interpretable, and defensible enough for the benchmark itself to support an honest first draft.

The discussion should be read with that scope in mind. This draft does not claim broad scene generalization, real-phone transfer, replay re-solving, or a final benchmark campaign. It is a focused second-pass closure attempt: one repaired fused runtime policy, one tuned fused nominal configuration, one compact three-condition suite, one artifact-driven paper build, and one local presentation bundle that makes the story inspectable by other people.

\section{Limitations}
Several limitations remain explicit. The suite still uses one scene and one robot preset. Replay remains analysis-only. The intermittent-anchor condition is a deterministic gating schedule rather than a photorealistic occlusion benchmark. The servo-stress condition is intentionally narrow and should be treated as a first draft of the actuation-realism story rather than a comprehensive stress campaign.

Those limits are acceptable for this draft because the goal is now narrower than a broad fused-advantage paper. This second-pass draft closes the minimum publishable stabilization question: whether the anchored visual--inertial estimator can be repaired enough to remain competitive and well calibrated on the nominal benchmark while avoiding catastrophic failure under intermittent anchor loss. It does not establish that fused estimation is generally superior to visual-only control in this setup, and the current suite still leaves visual as the stronger waypoint-tracking baseline.

\subsection{Threats To Validity}
The present stabilization result is intentionally narrow. It is demonstrated in one scene with one robot preset, deterministic anchor-suppression windows, a synthetic IMU, and auxiliary tags disabled in the final fused lock. Those choices were deliberate because they isolate the repaired fused runtime behavior, but they also limit how broadly the current draft should be generalized. The benchmark therefore supports an honest repair-and-stabilization claim, not a broad statement that fused estimation is now universally superior in anchored closed-loop control.

\bibliographystyle{plain}
\bibliography{references}

\end{document}
